\documentclass[aspectratio=169,14pt]{beamer}

% ─── Кодировка и язык ────────────────────────────────────────────────────────
\usepackage[utf8]{inputenc}
\usepackage[T2A]{fontenc}
\usepackage[russian]{babel}

% ─── Тема оформления ─────────────────────────────────────────────────────────
\usetheme{Madrid}
\usecolortheme{seahorse}

% Цветовая палитра (Ocean Gradient)
\definecolor{primary}{HTML}{065A82}
\definecolor{secondary}{HTML}{1C7293}
\definecolor{accent}{HTML}{02C39A}
\definecolor{lightbg}{HTML}{EAF4F8}
\definecolor{codebg}{HTML}{F0F4F8}

\setbeamercolor{palette primary}{bg=primary,fg=white}
\setbeamercolor{palette secondary}{bg=secondary,fg=white}
\setbeamercolor{palette tertiary}{bg=primary!80,fg=white}
\setbeamercolor{palette quaternary}{bg=primary,fg=white}
\setbeamercolor{structure}{fg=primary}
\setbeamercolor{block title}{bg=primary,fg=white}
\setbeamercolor{block body}{bg=lightbg}
\setbeamercolor{alerted text}{fg=accent!70!black}
\setbeamercolor{frametitle}{bg=primary,fg=white}

\setbeamerfont{title}{size=\large,series=\bfseries}
\setbeamerfont{frametitle}{size=\normalsize,series=\bfseries}

% Убираем навигационные символы
\setbeamertemplate{navigation symbols}{}

% Нумерация слайдов в колонтитуле
\setbeamertemplate{footline}{%
	\leavevmode%
	\hbox{%
		\begin{beamercolorbox}[wd=.45\paperwidth,ht=2.25ex,dp=1ex,left,leftskip=1em]{palette primary}%
			\usebeamerfont{author in head/foot}ЮФУ · Физический факультет · 2026
		\end{beamercolorbox}%
		\begin{beamercolorbox}[wd=.40\paperwidth,ht=2.25ex,dp=1ex,center]{palette secondary}%
			\usebeamerfont{title in head/foot}Администрирование ОС Linux
		\end{beamercolorbox}%
		\begin{beamercolorbox}[wd=.15\paperwidth,ht=2.25ex,dp=1ex,right,rightskip=1em]{palette primary}%
			\usebeamerfont{date in head/foot}\insertframenumber{} / \inserttotalframenumber
	\end{beamercolorbox}}%
}

% ─── Листинги кода ───────────────────────────────────────────────────────────
\usepackage{listings}
\usepackage{xcolor}

% Явное определение синтаксиса YAML для пакета listings
\lstdefinelanguage{yaml}{
	keywords={true,false,null,y,n,services,volumes,networks,driver,driver_opts,image,privileged,environment,networks,healthcheck,test,interval,timeout,retries,depends_on,condition},
	keywordstyle=\color{primary}\bfseries,
	basicstyle=\ttfamily\scriptsize,
	sensitive=false,
	comment=[l]{\#},
	commentstyle=\color{accent!70!black}\ttfamily,
	stringstyle=\color{secondary}\ttfamily,
	moredelim=[l][\color{primary}]{\:},
	morestring=[b]',
	morestring=[b]"
}

\lstset{
	backgroundcolor=\color{codebg},
	basicstyle=\ttfamily\scriptsize,
	breaklines=true,
	frame=single,
	rulecolor=\color{primary!30},
	numbers=none,
	showstringspaces=false,
	tabsize=2,
	extendedchars=true,
	inputencoding=utf8,
	literate=
	{а}{{\cyrа}}1 {б}{{\cyrb}}1 {в}{{\cyrv}}1 {г}{{\cyrg}}1 {д}{{\cyrd}}1
	{е}{{\cyre}}1 {ё}{{\cyryo}}1 {ж}{{\cyrzh}}1 {з}{{\cyrz}}1 {и}{{\cyri}}1
	{й}{{\cyrishrt}}1 {к}{{\cyrk}}1 {л}{{\cyrl}}1 {м}{{\cyrm}}1 {н}{{\cyrn}}1
	{о}{{\cyro}}1 {п}{{\cyrp}}1 {р}{{\cyrr}}1 {с}{{\cyrs}}1 {т}{{\cyrt}}1
	{у}{{\cyru}}1 {ф}{{\cyrf}}1 {х}{{\cyrh}}1 {ц}{{\cyrc}}1 {ч}{{\cyrch}}1
	{ш}{{\cyrsh}}1 {щ}{{\cyrshch}}1 {ъ}{{\cyrhrdsn}}1 {ы}{{\cyrery}}1
	{ь}{{\cyrsftsn}}1 {э}{{\cyrerev}}1 {ю}{{\cyryu}}1 {я}{{\cyrya}}1
	{А}{{\CYRA}}1 {Б}{{\CYRB}}1 {В}{{\CYRV}}1 {Г}{{\CYRG}}1 {Д}{{\CYRD}}1
	{Е}{{\CYRE}}1 {Ё}{{\CYRYO}}1 {Ж}{{\CYRZH}}1 {З}{{\CYRZ}}1 {И}{{\CYRI}}1
	{Й}{{\CYRISHRT}}1 {К}{{\CYRK}}1 {Л}{{\CYRL}}1 {М}{{\CYRM}}1 {Н}{{\CYRN}}1
	{О}{{\CYRO}}1 {П}{{\CYRP}}1 {Р}{{\CYRR}}1 {С}{{\CYRS}}1 {Т}{{\CYRT}}1
	{У}{{\CYRU}}1 {Ф}{{\CYRF}}1 {Х}{{\CYRH}}1 {Ц}{{\CYRC}}1 {Ч}{{\CYRCH}}1
	{Ш}{{\CYRSH}}1 {Щ}{{\CYRSHCH}}1 {Ъ}{{\CYRHRDSN}}1 {Ы}{{\CYRERY}}1
	{Ь}{{\CYRSFTSN}}1 {Э}{{\CYREREV}}1 {Ю}{{\CYRYU}}1 {Я}{{\CYRYA}}1
}

% ─── Доп. пакеты ─────────────────────────────────────────────────────────────
\usepackage{booktabs}
\usepackage{array}
\usepackage{tabularx}
\usepackage{tikz}
\usetikzlibrary{shapes.geometric,arrows.meta,fit,positioning,backgrounds}

% Безопасные строки авторов для PDF-закладок (убирает предупреждения)
\newcommand{\AuthorG}{\texorpdfstring{Григорян А.\,А.}{Григорян А.А.}}
\newcommand{\AuthorK}{\texorpdfstring{Карлашов А.\,В.}{Карлашов А.В.}}
\newcommand{\AuthorKr}{\texorpdfstring{Кроливец И.\,С.}{Кроливец И.С.}}

% ════════════════════════════════════════════════════════════════════════════════
%  МЕТАДАННЫЕ
% ════════════════════════════════════════════════════════════════════════════════
\title{Развёртывание и настройка NFS-сервера\\для сетевого хранения данных}
\subtitle{Курсовой проект по курсу «Администрирование Unix-подобных систем»}
\author{\AuthorG \and \AuthorK \and \AuthorKr}
\institute{ЮФУ · Физический факультет}
\date{Ростов-на-Дону, 2026}

% ════════════════════════════════════════════════════════════════════════════════
\begin{document}
	
	% ─── Титульный слайд ─────────────────────────────────────────────────────────
	\begin{frame}
		\titlepage
	\end{frame}
	
	% ─── Оглавление ──────────────────────────────────────────────────────────────
	\begin{frame}{Содержание}
		\tableofcontents[hideallsubsections]
	\end{frame}
	
	% ════════════════════════════════════════════════════════════════════════════════
	\section{Актуальность и постановка задачи}
	% ════════════════════════════════════════════════════════════════════════════════
	
	\begin{frame}{Актуальность проекта и проблематика}
		\begin{columns}[T]
			\begin{column}{0.6\textwidth}
				\begin{block}{Ключевые проблемы}
					\begin{itemize}
						\item \textbf{Изолированное хранение данных}~— каждый контейнер хранит данные отдельно, что порождает дублирование
						\item \textbf{Масштабируемость}~— требуется единая точка хранения для всех рабочих узлов
						\item \textbf{Configuration drift}~— ручная настройка приводит к расхождению конфигураций
					\end{itemize}
				\end{block}
			\end{column}
			\begin{column}{0.38\textwidth}
				\begin{alertblock}{Решение}
					\centering
					\textbf{NFSv4}\\[2pt]
					Единый порт 2049/TCP\\
					Kerberos-ready\\[6pt]
					\textbf{IaC подход}\\[2pt]
					Docker Compose\\
					bootstrap.sh\\
					CI/CD пайплайн
				\end{alertblock}
			\end{column}
		\end{columns}
	\end{frame}
	
	\begin{frame}{Цели и задачи работы}
		\begin{block}{\centering Цель работы}
			\centering\small
			Проектирование, развёртывание и защита NFS-сервера сетевого хранения данных в контейнеризированной среде под управлением ОС Linux
		\end{block}
		\vspace{0.2em}
		\textbf{Задачи для достижения цели:}
		\begin{enumerate}\small
			\item Развёртывание NFS-сервера в Docker с персистентным томом данных
			\item Подключение нескольких клиентских контейнеров через Docker Volume (\texttt{type: nfs4})
			\item Реализация сетевой изоляции (\texttt{storage-net}) и hardening (UFW, \texttt{cap\_drop})
			\item Написание скриптов тестирования (\texttt{test\_nfs.sh}) и бэкапа
			\item Настройка CI/CD пайплайна и мониторинга Prometheus + Grafana
		\end{enumerate}
	\end{frame}
	
	% ════════════════════════════════════════════════════════════════════════════════
	\section{Распределение ролей}
	% ════════════════════════════════════════════════════════════════════════════════
	
	\begin{frame}{Распределение ролей и индивидуальный вклад}
		\small
		\begin{columns}[T]
			\begin{column}{0.32\textwidth}
				\begin{block}{Григорян Ашот\\{\footnotesize Инженер по автоматизации}}
					\begin{itemize}\footnotesize
						\item Структура репозитория, Git
						\item Docker Compose манифест
						\item \texttt{bootstrap.sh}~— хост
						\item Dockerfile (non-root)
						\item Makefile: \texttt{up/down/test}
					\end{itemize}
				\end{block}
			\end{column}
			\begin{column}{0.32\textwidth}
				\begin{block}{Карлашов Андрей\\{\footnotesize Системный администратор / SRE}}
					\begin{itemize}\footnotesize
						\item Настройка \texttt{/etc/exports}
						\item NFSv4: fsid=0, порт 2049
						\item Клиентские тома (\texttt{type: nfs4})
						\item healthcheck: showmount
						\item HA-тест: failover
					\end{itemize}
				\end{block}
			\end{column}
			\begin{column}{0.32\textwidth}
				\begin{block}{Кроливец Иван\\{\footnotesize Инженер по безопасности}}
					\begin{itemize}\footnotesize
						\item \texttt{hardening.sh}: UFW
						\item Изоляция \texttt{storage-net}
						\item Prometheus + Grafana
						\item Дашборды: I/O, метрики
						\item CI/CD: GitHub Actions
					\end{itemize}
				\end{block}
			\end{column}
		\end{columns}
	\end{frame}
	
	% ════════════════════════════════════════════════════════════════════════════════
	\section{Теоретические основы}
	% ════════════════════════════════════════════════════════════════════════════════
	
	\begin{frame}{Теоретические основы: протокол NFS}
		\begin{columns}[T]
			\begin{column}{0.55\textwidth}
				\textbf{Ключевые свойства NFSv4:}
				\begin{itemize}\small
					\item Единственный порт \textbf{2049/TCP}~— упрощает правила брандмауэра
					\item Сессионная модель со стейтом
					\item Kerberos / RPCSEC\_GSS
					\item pNFS (v4.1)~— параллельный доступ
				\end{itemize}
				\vspace{0.3em}
				\textbf{Опции \texttt{/etc/exports}:}\\[2pt]
				\footnotesize\texttt{rw/ro}, \texttt{sync/async}, \texttt{root\_squash}, \texttt{no\_subtree\_check}, \texttt{fsid=0}
			\end{column}
			\begin{column}{0.43\textwidth}
				\begin{block}{Сравнение альтернатив}
					\small
					\begin{tabular}{ll}
						\toprule
						\textbf{Технология} & \textbf{Вывод} \\
						\midrule
						NFS   & \textcolor{accent!70!black}{\textbf{Оптимален}} \\
						Samba & Избыточен \\
						GlusterFS & Избыточен \\
						iSCSI & Другой класс \\
						\bottomrule
					\end{tabular}
				\end{block}
			\end{column}
		\end{columns}
	\end{frame}
	
	% ════════════════════════════════════════════════════════════════════════════════
	\section{Архитектура и сетевая инфраструктура}
	% ════════════════════════════════════════════════════════════════════════════════
	
	\begin{frame}{Архитектура распределённого кластера}
		\centering
		\begin{tikzpicture}[
			node distance=1.0cm and 1.4cm,
			box/.style={draw=primary,fill=lightbg,rounded corners=4pt,
				minimum width=2.4cm,minimum height=0.85cm,
				font= \small\bfseries,align=center,text=primary},
			net/.style={draw=secondary,fill=secondary!10,rounded corners=6pt,
				minimum width=9.5cm,minimum height=1.7cm,
				font=\footnotesize\itshape,align=left},
			arr/.style={-{Stealth[length=5pt]},thick,color=primary}
			]
			% storage-net background
			\begin{scope}
				\node[net,label={[font=\footnotesize\bfseries,color=secondary]above:storage-net (172.20.0.0/24)}]
				(snet) at (0,0) {};
			\end{scope}
			
			% nodes inside storage-net
			\node[box] (srv) at (-3.2,0)  {nfs-server\\{\tiny 172.20.0.10}};
			\node[box] (c1)  at (0,0)     {nfs-client-1\\{\tiny 172.20.0.11}};
			\node[box] (c2)  at (3.2,0)   {nfs-client-2\\{\tiny 172.20.0.12}};
			
			% monitoring
			\node[box,fill=accent!15,draw=accent!60!black] (mon) at (0,-2.0)
			{monitoring\\{\tiny Prometheus + Grafana}};
			
			% arrows
			\draw[arr] (srv) -- (c1);
			\draw[arr] (srv) -- (c2);
			\draw[arr,dashed,color=secondary] (srv) -- (mon) node[midway,right,font=\tiny]{monitoring-net};
			
			% volume label
			\node[font=\tiny,color=primary,below=2pt of srv] {nfs\_data (persistent)};
		\end{tikzpicture}
		
		\vspace{0.2em}
		\begin{columns}[T]
			\begin{column}{0.48\textwidth}
				\centering\footnotesize\textbf{storage-net}~— NFS-трафик изолирован, порт 2049 закрыт снаружи
			\end{column}
			\begin{column}{0.48\textwidth}
				\centering\footnotesize\textbf{monitoring-net}~— Grafana на порту 3000 с HTTP-auth
			\end{column}
		\end{columns}
	\end{frame}
	
	\begin{frame}[fragile]{Сетевая инфраструктура и изоляция}
		\begin{columns}[T]
			\begin{column}{0.55\textwidth}
				\begin{block}{Virtual Network \texttt{storage-net}}
					\begin{itemize}\footnotesize
						\item \texttt{nfs-server}: fixed IP \texttt{172.20.0.10}
						\item Клиенты: \texttt{.11} и \texttt{.12}
						\item Порт 2049 закрыт из внешнего мира
					\end{itemize}
				\end{block}
				\begin{block}{Мониторинг периметра}
					\begin{itemize}\footnotesize
						\item Grafana: порт 3000, HTTP-auth
						\item Prometheus: доступ только через localhost
					\end{itemize}
				\end{block}
			\end{column}
			\begin{column}{0.43\textwidth}
				\begin{alertblock}{Правила UFW}
					\scriptsize
					\begin{verbatim}
						ufw default deny incoming
						ufw allow 22/tcp
						ufw allow from \
						172.20.0.0/24 \
						to any port 2049
					\end{verbatim}
				\end{alertblock}
				\vspace{0.1em}
				\centering\footnotesize\textcolor{accent!70!black}{\textbf{Защита от утечек трафика}}\\
				NFS локализован внутри \texttt{storage-net}
			\end{column}
		\end{columns}
	\end{frame}
	
	% ════════════════════════════════════════════════════════════════════════════════
	\section{Автоматизация (IaC)}
	% ════════════════════════════════════════════════════════════════════════════════
	
	\begin{frame}[fragile]{Автоматизация развёртывания: Docker Compose}
		\begin{lstlisting}[language=yaml,basicstyle=\ttfamily\tiny]
			services:
			nfs-server:
			image: erichough/nfs-server:latest
			privileged: true
			environment:
			NFS_EXPORT_0: '/exports/shared *(rw,sync,no_subtree_check,fsid=0)'
			volumes:
			- nfs_data:/exports/shared
			networks:
			storage-net:
			ipv4_address: 172.20.0.10
			healthcheck:
			test: ["CMD", "showmount", "-e", "localhost"]
			interval: 30s  timeout: 10s  retries: 5
			
			nfs-client-1:
			image: alpine:3.19
			depends_on:
			nfs-server:
			condition: service_healthy
			volumes:
			- nfs_shared:/mnt/shared
			networks: [storage-net]
		\end{lstlisting}
	\end{frame}
	
	\begin{frame}[fragile]{Монтирование NFS через Docker Volume (NFSv4)}
		\begin{lstlisting}[language=yaml,basicstyle=\ttfamily\tiny]
			volumes:
			nfs_shared:
			driver: local
			driver_opts:
			type: nfs4
			o: "addr=172.20.0.10,rw,sync,nfsvers=4,hard,intr"
			device: ":/exports/shared"
			nfs_data:
			driver: local
			networks:
			storage-net:
			driver: bridge
			ipam:
			config:
			- subnet: 172.20.0.0/24
			monitoring-net:
			driver: bridge
		\end{lstlisting}
	\end{frame}
	
	\begin{frame}[fragile]{Первичная подготовка среды (\texttt{bootstrap.sh})}
		\begin{lstlisting}[language=bash,basicstyle=\ttfamily\tiny]
			#!/bin/bash
			# @file bootstrap.sh  @brief Host initialization
			set -euo pipefail
			
			EXPORT_DIR="${EXPORT_DIR:-/opt/nfs/shared}"
			BACKUP_DIR="${BACKUP_DIR:-/opt/nfs/backups}"
			
			# Initialization of target mount directories
			for dir in "${EXPORT_DIR}" "${BACKUP_DIR}"; do
			if [[ ! -d "${dir}" ]]; then
			mkdir -p "${dir}"
			chown nobody:nogroup "${dir}"
			chmod 777 "${dir}"
			echo "LOG: Created ${dir}"
			fi
			done
			
			# Environment config verification
			if [[ ! -f ".env" ]]; then
			cp .env.example .env
			echo "WARN: .env generated!"
			fi
		\end{lstlisting}
	\end{frame}
	
	% ════════════════════════════════════════════════════════════════════════════════
	\section{Безопасность}
	% ════════════════════════════════════════════════════════════════════════════════
	
	\begin{frame}{Обеспечение безопасности (Hardening)}
		\begin{columns}[T]
			\begin{column}{0.48\textwidth}
				\begin{block}{Уровень хоста (\texttt{hardening.sh})}
					\begin{itemize}
						\item UFW: \texttt{default deny incoming}
						\item 22/TCP~— SSH, только доверенные IP
						\item 2049/TCP~— строго для \texttt{172.20.0.0/24}
					\end{itemize}
				\end{block}
				\begin{block}{Уровень NFS}
					\footnotesize\texttt{root\_squash}: привилегированные запросы маппятся в \texttt{nobody} (UID 65534)
				\end{block}
			\end{column}
			\begin{column}{0.48\textwidth}
				\begin{block}{Уровень контейнеров}
					\begin{itemize}\footnotesize
						\item \texttt{cap\_drop: [ALL]} у клиентов
						\item \texttt{read\_only: true} для корня ФС
						\item Запуск процессов от \texttt{USER appuser}
						\item \texttt{privileged: true} для сервера~— необходимое исключение ядерного уровня
					\end{itemize}
				\end{block}
			\end{column}
		\end{columns}
	\end{frame}
	
	% ════════════════════════════════════════════════════════════════════════════════
	\section{Тестирование и CI/CD}
	% ════════════════════════════════════════════════════════════════════════════════
	
	\begin{frame}[fragile]{Интеграционное тестирование (\texttt{test\_nfs.sh})}
		\begin{columns}[T]
			\begin{column}{0.56\textwidth}
				\textbf{4 стадии автоматической валидации:}
				\begin{enumerate}\footnotesize
					\item \textbf{Сетевой скан}~— проверка экспорта через \texttt{showmount} (timeout 60s)
					\item \textbf{Монтирование}~— парсинг вывода \texttt{df -h} в клиентах
					\item \textbf{I/O валидация}~— сквозная запись из client-1 и чтение из client-2
					\item \textbf{HA-сценарий}~— рестарт сервера и проверка восстановления дескрипторов файлов
				\end{enumerate}
			\end{column}
			\begin{column}{0.42\textwidth}
				\begin{block}{Резервное копирование}
					\scriptsize
					\textbf{rsync} \texttt{--archive --checksum}\\[3pt]
					Инкрементальный бэкап (7 суток):
					\begin{lstlisting}[language=bash,basicstyle=\ttfamily\tiny]
						find "${BACKUP_DIR}" -mtime +7 -exec rm {} +
					\end{lstlisting}
				\end{block}
			\end{column}
		\end{columns}
	\end{frame}
	
	\begin{frame}{Конвейер CI/CD: GitHub Actions}
		\begin{columns}[T]
			\begin{column}{0.55\textwidth}
				\textbf{Спецификация \texttt{pr-validation.yml}}\\
				(Триггер: Pull Request в main):
				\begin{enumerate}\small
					\item \textbf{yamllint}~— валидация манифестов
					\item \textbf{ShellCheck}~— статический анализ bash
					\item \textbf{Hadolint}~— оптимизация Dockerfile
					\item \textbf{Инфратест}~— запуск \texttt{test\_nfs.sh}
				\end{enumerate}
			\end{column}
			\begin{column}{0.43\textwidth}
				\begin{alertblock}{Защита веток}
					\centering\footnotesize
					Слияние с веткой \texttt{main} заблокировано до прохождения тестов
				\end{alertblock}
				\vspace{0.1em}
				\centering\scriptsize
				\begin{tabular}{ll}
					\textcolor{accent!70!black}{\textbf{Анализ}} & 3 линтера \\
					\textcolor{accent!70!black}{\textbf{Тест}} & Стенд в облаке \\
				\end{tabular}
			\end{column}
		\end{columns}
	\end{frame}
	
	% ════════════════════════════════════════════════════════════════════════════════
	\section{Мониторинг}
	% ════════════════════════════════════════════════════════════════════════════════
	
	\begin{frame}{Мониторинг: Prometheus + Grafana}
		\begin{columns}[T]
			\begin{column}{0.52\textwidth}
				\textbf{Сбор телеметрии:}
				\begin{itemize}\footnotesize
					\item \texttt{node\_exporter}~— системные метрики хоста (CPU, RAM, Disk I/O)
					\item \texttt{cAdvisor}~— производительность контейнеров
					\item \textbf{Prometheus}~— TSDB, scrape: 15s
				\end{itemize}
			\end{column}
			\begin{column}{0.46\textwidth}
				\textbf{Grafana дашборды (IaC):}
				\begin{itemize}\footnotesize
					\item NFS метрики: ops/s, bytes r/w
					\item Ресурсы: утилизация лимитов
				\end{itemize}
				\begin{alertblock}{Алертинг}
					\centering\small
					Оповещение при дауне > 30с
				\end{alertblock}
			\end{column}
		\end{columns}
	\end{frame}
	
	% ════════════════════════════════════════════════════════════════════════════════
	\section{Итоги}
	% ════════════════════════════════════════════════════════════════════════════════
	
	\begin{frame}{Итоги проекта}
		\begin{center}
			\small
			\begin{tabular}{lc}
				\toprule
				\textbf{Результат} & \textbf{Статус} \\
				\midrule
				Развёртывание NFS Replica Set в Docker Compose & \textcolor{accent!70!black}{\checkmark\ Выполнено} \\
				Автоматизация: \texttt{bootstrap.sh} и Makefile & \textcolor{accent!70!black}{\checkmark\ Выполнено} \\
				Многоуровневый hardening (UFW, \texttt{cap\_drop}, RO) & \textcolor{accent!70!black}{\checkmark\ Выполнено} \\
				Интеграционные тесты и HA-тестирование & \textcolor{accent!70!black}{\checkmark\ Выполнено} \\
				Мониторинг Prometheus + Grafana & \textcolor{accent!70!black}{\checkmark\ Выполнено} \\
				CI/CD пайплайн GitHub Actions & \textcolor{accent!70!black}{\checkmark\ Выполнено} \\
				\bottomrule
			\end{tabular}
		\end{center}
		\vspace{-0.2em}
		\begin{block}{\centering Ключевые достижения}
			\centering\footnotesize
			Концепция \textbf{IaC} полностью исключила дрейф конфигураций \quad\textbullet\quad Сетевая безопасность.\\
			Наблюдаемость через Prometheus/Grafana \quad\textbullet\quad Контроль стабильности через CI/CD.
		\end{block}
	\end{frame}
	
	% ─── Финальный слайд ─────────────────────────────────────────────────────────
	\begin{frame}
		\begin{center}
			\vspace{0.5em}
			{\Large\bfseries\color{primary} Спасибо за внимание!}\\[1.0em]
			{\normalsize Григорян А.\,А. \quad Карлашов А.\,В. \quad Кроливец И.\,С.}\\[0.3em]
			{\small ЮФУ · Физический факультет · 2026}\\[0.8em]
			\begin{block}{\centering Тема проекта}
				\centering\small
				Развёртывание и настройка NFS-сервера для сетевого хранения данных
			\end{block}
			\vspace{0.3em}
			{\footnotesize Проверил: Тимошенко Павел Евгеньевич}
		\end{center}
	\end{frame}
	
\end{document}